\chapter{Conclusion and Recommendation}
\section{Conclusion}
Victor Miller and Neal Koblitz independently conceived the notion of utilizing elliptic curves as the basis of a group for the discrete logarithm problem in 1985, and this was the beginning of elliptic curve cryptography.The fact that no sub exponential technique existed to solve the discrete logarithm issue on a suitably designed elliptic curve is the major reason behind Elliptic curve cryptography's appeal. This meant that Elliptic curve cryptography used substantially lower parameters while maintaining the same level of security.Elliptic Curves provide security in the same way as traditional systems (such as RSA) do, but with fewer bits.\\
The use of elliptic curves in cryptography necessitated a smaller chip size, lower power consumption, and increased speed. In the contest of generating keys, the cyclic group, key pair generation together with the methods of application of rational numbers on elliptic curves in cryptography and finding points algebraically were used.\\
We also discussed about Diffie - Hellman Key Exchange on elliptic curve cryptography, ElGammal Public Key Encryption and ElGammal Digital Signatures which were geared towards the secure transfer of information. It was also found that, in channeling attacks and security twist the methods used were based on Attacks on Discrete logs, Baby Step, Giant Step, Pollard's p Method. A script was derived for generating private key cryptography.
 
 
 
\section{Recommendation}
Elliptic curve cryptography is considered to be a tool for mathematical studies, specifically in getting to know private key cryptography in ECC to keep confidential information between two parties, along with channel attacks and twist-security attacks based on the details and illustrative codes in this research.